% HD-F1-I03. Root proof; independent whole-proof review pending.
\section[HD-F1-I03: nine letters and two or three reflections]{HD-F1-I03: exact period ranks three and four for the rational nine-letter strata with two or three reflections}
\label{sec:f1-nine-letter-two-three-reflection-rank}

Retain the actual faithful monomial tower, primitive character,
polarization and period comparison of HD-I01--HD-I21.
For every component with
\[
 n=5,\qquad m=9,\qquad g(Y)=0,\qquad \nu=4,
 \qquad A=b-\nu\in\{2,3\},
\]
the original period differential has rank $A+1$ at every smooth
Hurwitz point. Thus the image dimensions are respectively three
and four. The parameters $d\in\{3,4,6\}$ and $e\mid d$ retain
their original meaning. Actual full $G(3,1,9)$ tuples below
establish nonemptiness of both strata.

\paragraph{1. The actual positive evaluation subbundle.}
Let $Q$ be the reduced divisor of uncolored cycles over the
four nonreflection values $B$, with every unramified uncolored
cycle retained. The I20 degree identity gives
\[
 \deg Q=(\nu-2)m+A-2n=8+A.
\]
Write $\mathcal V=h_*\mathcal O_Y(Q)$ on the actual target
$\mathbf P^1$. It is locally free of rank nine. For every
scheme-theoretic fibre $D_z=h^*(z)$, including a nonreduced
fibre, $\mathcal O_Y(Q-D_z)$ has degree $A-1$, so its
$H^1$ vanishes. The exact fibre restriction therefore makes
the evaluation $H^0(\mathcal V)\otimes\mathcal O\to\mathcal V$
surjective at every point. Moreover
\[
 h^0(\mathcal V)=9+A,\qquad \deg\mathcal V=A,\qquad
 U=H^0(\mathcal V(-1)),\quad \dim U=A.
\]
The degree formula follows from finite pushforward and Euler
characteristic; the section counts are those of the actual
line bundles on the rational source curve.

Consider the evaluation morphism, tensored by $\mathcal O(1)$,
\[
       e_+:U\otimes\mathcal O(1)\longrightarrow\mathcal V.
\]
Its fibre kernel at any $z$ is
$H^0(\mathcal V(-1)(-z))$. By projection formula this is
the section space of $\mathcal O_Y(Q)\otimes h^*\mathcal O(-2)$,
of degree $A-10<0$. It vanishes. Thus $e_+$ is fibrewise
injective. Locally a nonzero maximal minor gives an invertible
minor and elementary row operations identify its image with
a free direct summand. Its actual cokernel $\mathcal F$ is
therefore a vector bundle, and
\[
 0\longrightarrow U\otimes\mathcal O(1)
 \xrightarrow{e_+}\mathcal V\longrightarrow\mathcal F
 \longrightarrow0,
 \qquad \operatorname{rank}\mathcal F=9-A,\quad
 \deg\mathcal F=0.
\]
It is globally generated, as a quotient of the globally
generated $\mathcal V$. Since $H^1(\mathcal O(1))=0$,
the exact section sequence gives $h^0(\mathcal F)=9-A$.
Its evaluation is consequently a surjective map between
bundles of the same rank and is an isomorphism.
Lift a basis of $H^0(\mathcal F)$ to $H^0(\mathcal V)$.
Those actual lifted sections split this exact sequence,
giving the explicit evaluation-based isomorphism
\[
       \mathcal V\simeq
          \mathcal O(1)^{\oplus A}\oplus
          \mathcal O^{\oplus(9-A)}.                 \tag{F1I03.1}
\]
The first summand is precisely the image of $e_+$.

\paragraph{2. The prescribed unit has a nonzero constant component.}
The literal function $1$ defines the actual unit section of
$\mathcal V$. Write it in (F1I03.1) as $(s,c)$.
We prove $c\ne0$ using the original map $h$ and divisor $Q$.

Choose homogeneous coordinates on the actual rational source
and target. The pullbacks of target coordinates are coprime
degree-nine polynomials $p_0,p_1$, so $h=[p_0:p_1]$.
Choose an isomorphism $\mathcal O_Y(Q)\simeq\mathcal O_Y(8+A)$.
It carries the section $1$ to a nonzero polynomial $q_Q$
whose zero divisor is exactly the reduced $Q$.
The induced isomorphism
$\mathcal O_Y(Q)\otimes h^*\mathcal O(-1)
 \simeq\mathcal O_Y(A-1)$ identifies the section map of
$e_+$ with the actual multiplication map
\[
 H^0(\mathcal O_Y(A-1))\otimes H^0(\mathcal O_{\mathbf P^1}(1))
 \longrightarrow H^0(\mathcal O_Y(8+A)),
 \quad (\alpha,\beta)\longmapsto\alpha p_0+\beta p_1.
\]
Consequently $c=0$ would give actual homogeneous polynomials
$\alpha,\beta$ of degree $A-1$, not both zero, satisfying
\[
                    q_Q=\alpha p_0+\beta p_1.       \tag{F1I03.2}
\]
Take their homogeneous greatest common factor $g$, of degree
$c_0$, and write $\alpha=g\alpha_0$, $\beta=g\beta_0$.
The coprime pair defines the actual morphism
\[
          \phi=[-\beta_0:\alpha_0]:Y\longrightarrow\mathbf P^1.
\]

When $\phi$ is nonconstant, its degree is
$e=A-1-c_0$. At every point of $Q$ outside the zero divisor
of $g$, equation (F1I03.2) gives $h=\phi$, because both
coordinate pairs have no simultaneous zero. Since $h(Q)$
is contained in the four-point set $B$, these points of $Q$
belong to $\phi^{-1}(B)$. A degree-$e$ map has at most
$4e$ distinct preimages of those four values. The divisor
$Q$ is reduced, so
\[
       8+A=\deg Q\le c_0+4e
          =4A-4-3c_0\le4A-4.                       \tag{F1I03.3}
\]
For $A=2$ this says $10\le4$, and for $A=3$ it says
$11\le8$. Both are impossible.

When $\phi$ is constant, let its value be $b_0$.
Outside the zeros of $g$, every point of $Q$ lies in the
single actual fibre $h^{-1}(b_0)$. If $b_0\notin B$,
there are no such points, contrary to $\deg Q>\deg g$.
If $b_0\in B$, at most eight reduced uncolored points lie
over it: nine would be nine uncolored cycles of total
length nine, hence nine unramified uncolored cycles, making
that actual monomial inertia the identity. All original
branch inertias are nonidentity. Thus
\[
             \deg Q\le8+\deg g\le8+(A-1)<8+A,
\]
again impossible. This includes the case where one of
$\alpha,\beta$ is zero. It proves $c\ne0$ in every case.

\paragraph{3. The actual quotient and original period differential.}
A constant automorphism of $\mathcal O^{9-A}$ sends $c$
to its first basis vector. The triangular automorphism
$(u,w)\mapsto(u-s w_1,w)$ then sends the unit to
$(0,1,0,\ldots,0)$. Quotienting by this specified unit gives
\[
       \mathcal W_Q=\mathcal V/\mathcal O
          \simeq\mathcal O(1)^{\oplus A}
                 \oplus\mathcal O^{\oplus(8-A)}.    \tag{F1I03.4}
\]
The exact I21 tangent-fibre morphism therefore identifies
\[
 \ker d\mathcal F\simeq H^0(\mathcal W_Q(-2))
       =H^0\bigl(\mathcal O(-1)^{\oplus A}
                    \oplus\mathcal O(-2)^{\oplus(8-A)}\bigr)=0.
\]
The source dimension is $b-3=A+1$. Both primitive
eigendifferential bundles have degree four on $Y$, so their
actual multiplication is onto $H^0(\mathcal O_Y(8))$:
every degree-eight monomial is a product of two degree-four
monomials. Compose with the original quadratic trace of I20
and the polarized isogeny of I02. That isogeny is the
restriction of the actual $q^*$ with
$q_*q^*=[|\ker\chi|]$ and the retained theta multiplier.
These are precisely the maps defining the original period
codifferential. It has the same rank as $d\mathcal F$,
namely $A+1$, at every smooth point of each stated component.

\paragraph{4. Actual full-group tuples in both strata.}
Use the actual $C,D\in G(3,1,9)$ of HD-F1-I01, with
$C$ having permutation $(1\,2\,3\,4\,5\,6)$, zero color
on this cycle, and color one on fixed letters $7,8,9$.
The permutation of $D$ is $(6\,7\,8\,9)$ with cycle
color one; its fixed letters $1,2$ have colors one and
two, and fixed letters $3,4,5$ have color zero.
For $a\in\mathbf Z/3$ let $R_i(a)$ be the diagonal matrix
with entry $\zeta_3^a$ at coordinate $i$ and one elsewhere.
Put
\[
 D''=D R_1(1)^{-1}R_2(2)^{-1},\qquad
 C'=C R_7(1)^{-1}.
\]
The first tuple is
\[
       (C,C^{-1},D'',R_1(1),R_2(2),D^{-1}),          \tag{F1I03.5}
\]
and the second is
\[
       (C',R_7(1),C^{-1},D'',R_1(1),R_2(2),D^{-1}). \tag{F1I03.6}
\]
Their products are one, using the literal identities
$D''R_1(1)R_2(2)=D$ and $C'R_7(1)=C$.
They therefore generate the original full group
$\langle C,D\rangle=G(3,1,9)$ proved in HD-F1-I01.
Their moving-rank lists are respectively
\[
           (8,8,4,1,1,6),\qquad (7,1,8,4,1,1,6).
\]
Thus precisely two and three entries are reflections,
and each tuple has exactly four nonreflection values.
The permutation ramification is unchanged at
$5+5+3+3=16$. Riemann--Hurwitz consequently gives
genus zero for the degree-nine intermediate quotient.

For (F1I03.5) and (F1I03.6), the colored-cycle counts are
respectively
\[
 (3,3,1,1,1,3),\qquad (2,1,3,1,1,1,3).
\]
In both cases the total is twelve.
Each removed fixed-cycle color reappears exactly as the
color of its new diagonal reflection. Thus both primitive
fractional-color sums remain six, and both holomorphic
multiplicities are $-1+6=5$.
The reduced uncolored counts over the four nonreflection
values are respectively $(1,1,5,3)$ and $(2,1,5,3)$,
giving $\deg Q=10,11$ as used above.
Riemann existence applied to these actual full-group
product-one tuples gives connected covers and their
specified monomial quotient towers. This establishes
both nonempty strata with all stated invariants.

\paragraph{Framework and exact remaining scope.}
All new morphisms are constructed on the original tower:
the positive evaluation subbundle, its lifted constant
quotient, the polynomial form of the prescribed unit,
the degree-bounded morphism $\phi$, and the triangular
unit-preserving quotient. The original local system,
Jacobian corner, polarized isogeny, quadratic trace,
Hilbert and cycle realizations, and normalized Poincare
comparison retain their F01--F04 maps. The coherent
exact sequences above give their actual $K_0$ relations;
Adams operations apply to those same classes.

This completes the two individual component strata of
T4-INDIVIDUAL. It retains all other values of $m,g,\nu,A$
in the exact I21 quotient-bundle and multiplication-image
gate. The full map and unit argument is proved here;
no generic splitting type is assigned to another stratum.
The published nonvanishing comparison with Filom,
\emph{On the Rank of the Thurston Pullback Map},
IMRN2023, Theorem1.3, remains the exact marked-Hurwitz
comparison of HD-F1-I01--I02: there are twelve colored
source points, degree nine, and respectively six or seven
marked target branch values. Its first assertion gives
nonvanishing because $12>9+2$. The actual evaluation
and unit morphisms above compute the full ranks three
and four. No priority claim for the nonvanishing is made.
